{\large\textbf{Acoustic levitation (20 pts)}}

An acoustic levitator consists of transducers which convert an electrical
signal into intense sound waves. The acoustic pressure of the resulting
interference pattern can levitate small objects close to the nodes of the
pattern.

\textbf{Warning:} Avoid short-circuiting the connections on the top and bottom
of the levitator.

The tasks E1-E7 of this exam are designed to be independent of each other's
results. You may skip tasks you are struggling with.

For the following tasks, consider only the nodes on the z-axis. For task E3 and
onward, use only the centremost node. The Mode X switch on the electrical box
(F), see Fig. 1, should be in position O.

\begin{center}\includegraphics[width=0.39\linewidth]{figures/EuPhO_2026_Q4_fig1.jpg}\end{center}

\textbf{Task E1:\quad Probing the acoustic field (1.0 pts)}

To get familiar with the acoustic field, use the green LEDs (A) and the
transducer (B) to construct a tool that emits light according to the local
acoustic intensity. A transducer is a capacitor whose charge depends on the
mechanical pressure applied to it.

a) Draw a circuit diagram for your tool. (0.5 pts)

b) Use your tool to roughly estimate the distance between two adjacent
antinodes on the $z$ axis in the centre of the levitator. (0.5 pts)

\textbf{Task E2:\quad Acoustic levitator properties (4.0 pts)}

\textit{In what follows, to levitate a small solid object, the voltage of the
connected power supply can be initially put to a high voltage, for example
$17\,\mathrm{V}$. The object can be placed in the nodes by using the tweezers.
Note that the levitating object can oscillate horizontally due to
air-turbulence; wait a little for things to settle. Small oscillations are not
detrimental to measurement results.}

a) Design, sketch and use a setup to determine as precisely as possible the
frequency of the transducers. The speed of sound in air is
$v \approx 340\,\mathrm{m/s}$. (2.0 pts)

% NOTE: the switch position after "to position" is drawn as a bold vertical bar (not in the text layer); reproduced as a bold sans I
b) Switch \textit{Mode X} to position \textsf{\textbf{I}} on the electrical box
(F) to change to a different mode of operation where the transducers in the
upper half of the levitator are supplied with a \textit{slightly} different
frequency $(f + \Delta f)$ than the lower one. Describe the changed behaviour
of the levitated objects in this mode and determine the frequency difference
$\Delta f$ and its sign. (2.0 pts)

\textbf{Remember to return the \textit{Mode X} switch to position O for the
rest of the exam.}

\textbf{Task E3:\quad Density of solid beads (3.0 pts)}

Levitated objects will be situated slightly below the node due to gravity. The
restoring acoustic force $\vec{F}$ acting on the levitated object is
proportional to $P^2$, where $P \propto U$ is the acoustic pressure amplitude
at the antinodes and $U$ is the voltage from the power source.

Experimentally determine the density of the unknown beads (O) as precisely as
possible. The density of the glass beads (P) is
$\rho_\mathrm{glass} = 2500\,\mathrm{kg/m^3}$. These spherical beads (O, P)
have the same size. (3.0 pts)

\textit{Hint:} Depending on your approach, for spherical objects and small
displacements $\Delta\vec{z}$ from the node, you may use the approximation:
\begin{equation}
\vec{F} = -f(R)P^2\Delta\vec{z} \tag{1}
\end{equation}
Here, $f(R)$ is some unknown function of the radius $R$ of the levitated
object.

\textbf{Task E4:\quad Evaporation (3.0 pts)}

\textit{For the remaining tasks, you will levitate liquids, which will form
small droplets in the nodes. Place droplets in the levitator using the
syringes (R) with needles (S). To avoid cross-contamination, use a separate
syringe for each liquid. A good starting voltage for most droplet sizes is
8–12 V; at higher voltages the droplets may explode.}

\textbf{Warning: Avoid spilling liquids on the electronics and the
transducers.} Keep the metal net (L) on the bottom of the levitator; if it gets
wet, remove it and let it dry. If the transducers seem to malfunction, raise
your HELP sign.

The exact shape of the levitated droplets can be complicated. Moderately
deformed droplets may be approximated as an ellipsoid with semi-axes $a, b, c$,
which has the volume:
\begin{equation}
V = \frac{4\pi}{3}abc \tag{2}
\end{equation}

In this task, you will study the evaporation of a liquid droplet in the
levitator. A simplified description for the evaporation of a droplet with
equivalent diameter $D \gtrsim 1.5\,\mathrm{mm}$ is given by:
\begin{equation}
\frac{d(D^2)}{dt} = -\gamma \tag{3}
\end{equation}
where $D$ is the equivalent diameter of the droplet, i.e., the diameter of a
sphere with the same volume $V$ as the real droplet, and $\gamma$ is the
evaporation constant.

\textbf{For this task, consider only liquid I and keep the voltage constant.}

a) Determine the evaporation constant $\gamma$ in units $\mathrm{m^2/s}$ for
droplets with an equivalent diameter $D \gtrsim 1.5\,\mathrm{mm}$. (2.5 pts)

b) Would you over- or underestimate the time it takes for a droplet to
evaporate entirely, using your result from a)? (0.5 pts)

\textbf{Task E5:\quad Surface tension (2.5 pts)}

The shape of the droplet depends only on its volume $V$ and the ratio
$P^2/\sigma$. Here, $\sigma$ is the surface tension of the liquid.

a) Perform a measurement to find out which liquid, I or II, has a lower
surface tension. (0.5 pts)

b) Experimentally determine the surface tension of liquid I, given that the
surface tension of liquid II is $0.073\,\mathrm{N/m}$. (2.0 pts)

\textbf{Task E6:\quad Explosion (4.0 pts)}

When the acoustic pressure is too high, the levitating droplet will explode. A
theoretical formula for the maximum voltage that can be applied to the droplet
before it explodes is:
\begin{equation}
U_\mathrm{max} = \sqrt{\frac{\alpha}{D} + \beta} \tag{4}
\end{equation}
where $\alpha$ and $\beta$ are some constants.

\textbf{Consider only liquid II for this task.}

a) Determine the constants $\alpha$ and $\beta$ in the units
V\textsuperscript{2}·mm and V\textsuperscript{2}, respectively. (2.0 pts)

b) Estimate the maximum diameter of a droplet that can be levitated. (2.0 pts)

\textbf{Task E7:\quad Mysterious line (2.5 pts)}

When the levitating droplet becomes rather flat, a levitator voltage can be
found for which the light from the red LED passing through the droplet forms a
sharp horizontal line on the screen behind.

\textbf{Consider only liquid III for this task.}

Draw a sketch of the ray propagation inside the droplet when the line appears.
Use the line phenomenon described above to estimate the refractive index $n$
of liquid III. (2.5 pts)

\textit{Hint:} For the line to be visible, the LED must be well aligned with
the droplet, see Fig. 2.

\textbf{Task E8:\quad Intentional damage penalty (-0.5 pts)}

All equipment and packaging is left on the table after the competition in
working condition. For breaking components or packaging due to careless
handling or failing to follow the warning in red text, a penalty will be
applied. (-0.5 pts)

\textbf{Equipment}

A.\quad Two green Light Emitting Diodes (LEDs). Do not connect the LEDs to the
power supply (G).

B.\quad Individual transducer.

C.\quad Metal frame with a levitator module and a red LED, which can be modeled
as a \textit{point source}.

D.\quad Capture unit for collecting falling objects.

E.\quad 2 hex keys for mounting and aligning the setup.

F.\quad Electrical enclosure with switches \textit{LED} and \textit{Mode X}.

G.\quad Power supply.

H.\quad Banana cables.

I.\quad Screen.

J.\quad Tape.

K.\quad Transparent ruler.

L.\quad Metal net.

M.\quad Tweezers.

N.\quad 15 blue styrofoam beads.

O.\quad Multiple white spherical beads of unknown material
($\varnothing\ 2.0\,\mathrm{mm}$).

P.\quad Multiple transparent spherical glass beads
($\varnothing\ 2.0\,\mathrm{mm}$, $\rho_\mathrm{glass} = 2500\,\mathrm{kg\,m^{-3}}$).

Q.\quad 3 different unknown liquids labelled I, II and III. They are all safe
to touch with your hands, but should not be ingested.

R.\quad 3 syringes.

S.\quad 3 blunt end needles for the syringes (be careful!).

T.\quad 2 clothes pegs.

U.\quad Stop watch.

V.\quad Lamp.

W.\quad Power cord.

X.\quad Piece of paper to dry the metal net (not in image).

If you suspect that a component has failed for some reason during the
competition, inform the invigilator to get help.

\begin{center}\includegraphics[width=1.00\linewidth]{figures/EuPhO_2026_Q4_fig2.jpg}\end{center}

\begin{center}Figure 1: Overview of the equipment.\end{center}

\begin{center}\includegraphics[width=0.53\linewidth]{figures/EuPhO_2026_Q4_fig3.jpg}\end{center}

Figure 2: Alignment of LED module. You can align the setup by using the screws
shown in the figure. Arrows labelled \textbf{1} are for horizontal movement and
tilting of the LED module, and arrows labelled \textbf{2} are for vertical
movement and tilting of the LED module.
